\begin{abstract}
Evaluating defenses against Advanced Persistent Threats (APTs) is difficult because existing approaches often measure isolated detections, attack techniques, or product-specific outcomes without preserving how individual defense decisions affect a multi-stage attack campaign. As a result, a high-level detection rate or aggregate score may reveal whether a defense performs well, but provides limited guidance on where the defensive chain failed and what should be repaired.

We present an APT defense evaluation framework that connects playbook-faithful adversary execution with action-level defense observations and repair-oriented scoring. The framework models each attack occurrence as a distinct experimental unit and binds its observed defense outcome to a scoring opportunity representing the first failed victim-side control point. It further incorporates canonical attack phases, defense-technology mappings, and stage-level dependency semantics to aggregate action outcomes into stage, playbook, and overall benchmark scores while preserving a reverse trace from each score to the underlying attack action and control failure.

We construct the benchmark from 53 multi-stage APT playbooks containing 1,796 authoritative raw action occurrences and design an execution protocol that retains the complete action population rather than selecting only representative techniques. Actions that cannot be reproduced natively are explicitly distinguished from exact replay and must preserve attack semantics and defender-relevant telemetry under controlled substitution. The scoring model distinguishes prevention, blocking, detection, and missed actions and supports capability-specific analysis in addition to the overall score.

Our current evaluation executes the benchmark in a controlled networked environment against reproduced defense systems. [FINAL EXPERIMENTAL RESULT SENTENCE: summarize replay coverage, evaluated systems, principal score differences, and one or two repair-oriented findings after the real experiments are complete.] These results are intended to show that APT defense evaluation should measure not only whether attacks are detected, but also where the attack chain remains viable and which defensive control should be strengthened first.
\end{abstract}